% Options for packages loaded elsewhere
\PassOptionsToPackage{unicode}{hyperref}
\PassOptionsToPackage{hyphens}{url}
\PassOptionsToPackage{dvipsnames,svgnames,x11names}{xcolor}
\documentclass[
  10pt,
]{article}
\usepackage{xcolor}
\usepackage[margin=0.9in]{geometry}
\usepackage{amsmath,amssymb}
\setcounter{secnumdepth}{-\maxdimen} % remove section numbering
\usepackage{iftex}
\ifPDFTeX
  \usepackage[T1]{fontenc}
  \usepackage[utf8]{inputenc}
  \usepackage{textcomp} % provide euro and other symbols
\else % if luatex or xetex
  \usepackage{unicode-math} % this also loads fontspec
  \defaultfontfeatures{Scale=MatchLowercase}
  \defaultfontfeatures[\rmfamily]{Ligatures=TeX,Scale=1}
\fi
\usepackage{lmodern}
\ifPDFTeX\else
  % xetex/luatex font selection
\fi
% Use upquote if available, for straight quotes in verbatim environments
\IfFileExists{upquote.sty}{\usepackage{upquote}}{}
\IfFileExists{microtype.sty}{% use microtype if available
  \usepackage[]{microtype}
  \UseMicrotypeSet[protrusion]{basicmath} % disable protrusion for tt fonts
}{}
\makeatletter
\@ifundefined{KOMAClassName}{% if non-KOMA class
  \IfFileExists{parskip.sty}{%
    \usepackage{parskip}
  }{% else
    \setlength{\parindent}{0pt}
    \setlength{\parskip}{6pt plus 2pt minus 1pt}}
}{% if KOMA class
  \KOMAoptions{parskip=half}}
\makeatother
\usepackage{longtable,booktabs,array}
\usepackage{calc} % for calculating minipage widths
% Correct order of tables after \paragraph or \subparagraph
\usepackage{etoolbox}
\makeatletter
\patchcmd\longtable{\par}{\if@noskipsec\mbox{}\fi\par}{}{}
\makeatother
% Allow footnotes in longtable head/foot
\IfFileExists{footnotehyper.sty}{\usepackage{footnotehyper}}{\usepackage{footnote}}
\makesavenoteenv{longtable}
\setlength{\emergencystretch}{3em} % prevent overfull lines
\providecommand{\tightlist}{%
  \setlength{\itemsep}{0pt}\setlength{\parskip}{0pt}}
\usepackage{bookmark}
\IfFileExists{xurl.sty}{\usepackage{xurl}}{} % add URL line breaks if available
\urlstyle{same}
\hypersetup{
  colorlinks=true,
  linkcolor={blue},
  filecolor={Maroon},
  citecolor={Blue},
  urlcolor={blue},
  pdfcreator={LaTeX via pandoc}}

\author{}
\date{}

\begin{document}

\section{Project Proposal (Week 2) --- Fire Detection \& Autonomous
Emergency
Response}\label{project-proposal-week-2-fire-detection-autonomous-emergency-response}

D7065E, Week 2 Proposal

\subsection{Summary}\label{summary}

\begin{longtable}[]{@{}
  >{\raggedright\arraybackslash}p{(\linewidth - 2\tabcolsep) * \real{0.5000}}
  >{\raggedright\arraybackslash}p{(\linewidth - 2\tabcolsep) * \real{0.5000}}@{}}
\toprule\noalign{}
\endhead
\bottomrule\noalign{}
\endlastfoot
\textbf{Team} & Nagarajan Ganesan (naggan-4) \\
\textbf{Use case} & Fire detection and autonomous emergency response \\
\textbf{Pitch} & Detect real fires (not false alarms) from correlated
smoke/CO/temperature readings, and autonomously trigger sprinklers,
fire-door control, evacuation routing, and emergency-responder
notification. \\
\textbf{Repository} & https://github.com/ennc0d3-ltu/d7065e-project \\
\end{longtable}

\subsection{1. Use case, and why}\label{use-case-and-why}

This project detects fire in a BuildSim office floor and drives an
autonomous emergency response: sprinklers, fire-door control, and
evacuation guidance. It was chosen because it has a genuine tension ---
reacting fast to a real fire versus not triggering sprinklers and
lockdowns on cooking smoke or dust --- which makes the detection
decision interesting, and because BuildSim's room adjacency and walkable
graph give a natural substrate for both smoke spread and evacuation
routing.

\subsection{2. Scenario}\label{scenario}

Smoke concentration rises in one room. CO and temperature begin rising
shortly after in a correlated way. The system fuses the three signals,
confirms a real fire (rather than a cooking or dust false alarm), and
responds: sprinklers activate in the affected room, nearby fire doors
unlock, an evacuation route is computed away from the affected zone, and
a simulated call is placed to the building's emergency-response team
(fire brigade / security) with the affected room, severity, and a live
occupancy status --- who is still in the affected zone, who has reached
a safe area, and who is unaccounted for --- updated as people move. The
3D dashboard highlights the room red, the spread risk orange, and the
safe route green. As smoke and heat recede, sprinklers deactivate.

\subsection{3. What will be built (rough
sketch)}\label{what-will-be-built-rough-sketch}

\begin{itemize}
\tightlist
\item
  \textbf{Sensors:} smoke, CO, and temperature per room, registered with
  BuildSim.
\item
  \textbf{Physical simulator:} advances per-room smoke/CO/temperature
  and occupancy dynamics, reads actuator effects, and generates the
  underlying ``ground truth'' that sensors sample.
\item
  \textbf{Autonomous decision-maker:} fuses smoke/CO/temperature to
  decide real-fire vs. false alarm, and commands the response.
\item
  \textbf{Actuators:} sprinklers and fire-door locks, writing state back
  to BuildSim.
\item
  \textbf{Evacuation planner:} computes a route away from affected rooms
  using BuildSim's walkable graph, re-routing as the affected zone
  changes.
\item
  \textbf{Occupancy tracker:} reads BuildSim's occupancy/entity data per
  room to maintain a live status --- in danger zone, en route, safe,
  unaccounted-for --- feeding both the dashboard and the notifier.
\item
  \textbf{Emergency-response notifier:} on a confirmed fire, alerts a
  simulated emergency-response team (fire brigade / building security)
  with affected room, severity, and the current occupancy status, so
  responders know who still needs help and where; logs
  acknowledgement/response status.
\item
  \textbf{Data pipeline / storage:} collect sensor readings and store
  recent history in TimescaleDB (PostgreSQL time-series extension),
  feeding both the detector and the dashboard.
\item
  \textbf{Dashboard:} room highlights (red/orange/green), affected
  zones, and evacuation route overlaid on the BuildSim 3D view.
\end{itemize}

Data flow (first version): simulator -\textgreater{} sensors
-\textgreater{} BuildSim -\textgreater{} TimescaleDB -\textgreater{}
detector -\textgreater{} actuators -\textgreater{} simulator, with the
dashboard reading BuildSim state plus selected history.

Each major component (sensors, detector, actuators, evacuation planner,
occupancy tracker, notifier, storage, and dashboard) runs as a separate
process communicating through defined interfaces, and talks to BuildSim
through its API. Each process is independently restartable and
re-registers with BuildSim on restart; requests use timeouts and retries
so a slow or crashed component degrades the system rather than freezing
it. The occupancy tracker is the single owner of
\texttt{PUT\ /api/occupancy} --- since it replaces the whole collection,
exactly one process may write it, to avoid two writers erasing each
other's updates.

\subsection{4. The autonomous part}\label{the-autonomous-part}

It decides, per room, whether readings indicate a real fire and what
response to trigger. Plain thresholds alone are not sufficient (false
alarms from cooking/dust need to be rejected with real confidence), so
the data-driven component is required rather than optional here:
\textbf{starting approach} is a lightweight anomaly/fusion model ---
e.g.~a one-class or isolation-forest model trained on normal
multi-sensor (smoke/CO/temperature) patterns, so a real fire is flagged
as a joint deviation rather than any single threshold crossing. This
still runs as a fast/reflex decision: local, bounded, no external calls,
always available. Plain correlated thresholds are the fallback if
training data proves too thin, but the model-based approach is the
target for the first working loop, not an optional upgrade.

\subsection{5. Biggest unknowns \& risks}\label{biggest-unknowns-risks}

\begin{itemize}
\tightlist
\item
  Whether the simulated smoke-spread between adjacent rooms is realistic
  enough that detection and evacuation decisions mean anything.
\item
  Tuning thresholds so that correlated cooking/dust events don't trigger
  a false response, without missing a real fire.
\item
  Whether evacuation re-routing against BuildSim's walkable graph is
  fast enough to update as the affected zone grows.
\item
  What a realistic emergency-responder notification interface looks like
  when there is no real external service to call (a mock endpoint, or a
  logged/simulated dispatch).
\item
  Whether BuildSim's occupancy/entity data is granular and timely enough
  to reliably classify people as ``in danger,'' ``en route,'' or
  ``safe'' during a fast-moving scenario.
\end{itemize}

\subsection{6. Physical simulation
approach}\label{physical-simulation-approach}

BuildSim provides the rooms; the physics and the people are simulated as
follows.

\begin{itemize}
\tightlist
\item
  \textbf{Physical processes:} a simple per-room model --- smoke
  concentration rises with fire intensity and spreads to adjacent rooms
  through corridors/doorways using the room adjacency graph; temperature
  rises with fire intensity; CO rises correlated with smoke. Sprinkler
  activation reduces smoke concentration and temperature in the room,
  closing the feedback loop.
\item
  \textbf{People movement (occupancy):} a standard weekday schedule
  (arrivals 7--9am, lunch 11:30--13:00, departures 16:00--18:00, empty
  weekends), with some randomness, so occupancy data is available to
  inform evacuation routing. During an evacuation, people are moved
  along BuildSim's walkable graph toward the planned safe route, so
  their live room/position can be read back and used to classify each
  person's status for the occupancy tracker.
\item
  \textbf{Public dataset (baseline reference):} no dataset in the
  course-provided list is fire-specific, so smoke/CO/temperature
  dynamics are physics-modelled rather than replayed. The UCI/Kaggle
  \textbf{Occupancy Detection} dataset (light, temperature, humidity,
  CO2 with ground-truth occupancy) is used as a reference for realistic
  ``normal, no-fire'' sensor baselines and occupancy correlation, which
  the fire-vs-false-alarm model is trained against.
\end{itemize}

The models are intentionally simple, but actuator actions (sprinklers)
must influence future sensor readings so the control loop closes.

\subsection{7. First step \& rough plan}\label{first-step-rough-plan}

The first step is \textbf{one room working end-to-end}: smoke/CO/temp
sensors → detector → sprinkler actuator → BuildSim, to prove the
fire-vs-false-alarm loop closes and sprinklers visibly reduce
smoke/temperature. After that: smoke spread to adjacent rooms, fire-door
control, the evacuation planner, occupancy tracker, the
emergency-response notifier, storage, and the dashboard are added, with
failure handling and evaluation (fault injection, false-alarm testing)
following toward the demo.

\end{document}
